% Part: first-order-logic
% Chapter: tableaux
% Section: identity

\documentclass[../../../include/open-logic-section]{subfiles}

\begin{document}

\olfileid[tr]{fol}{tab}{ide}

\olsection{\usetoken{S}{identity} İçeren \usetoken{P}{tableau}}

!!^{tableau}s, !!{identity} içerdiklerinde ek çıkarım kuralları gerektirir.
$\eq$ kuralları şöyledir ($t$, $t_1$ ve $t_2$ kapalı terimlerdir):

\begin{defish}
\AxiomC{}
\RightLabel{$\eq$}
\UnaryInfC{\sFmla{\True}{\eq[t][t]}}
\DisplayProof
\hfill
\AxiomC{\sFmla{\True}{\eq[t_1][t_2]}}
\noLine
\UnaryInfC{\sFmla{\True}{!A(t_1)}}
\RightLabel{$\TRule{\True}{\eq}$}
\UnaryInfC{\sFmla{\True}{!A(t_2)}}
\DisplayProof
\hfill
\AxiomC{\sFmla{\True}{\eq[t_1][t_2]}}
\noLine
\UnaryInfC{\sFmla{\False}{!A(t_1)}}
\RightLabel{$\TRule{\False}{\eq}$}
\UnaryInfC{\sFmla{\False}{!A(t_2)}}
\DisplayProof
\end{defish}
Diğer bütün kurallardan farklı olarak $\TRule{\True}{\eq}$ ve
$\TRule{\False}{\eq}$ kurallarının dalda önceden \emph{iki} işaretli
!!{formula}s, yani hem $\sFmla{\True}{\eq[t_1][t_2]}$ hem de
$\sFmla{S}{!A(t_1)}$ bulunmasını gerektirdiğine dikkat edin.

\begin{ex}
$s$ ile $t$ kapalı terimlerse $\eq[s][t], !A(s) \Proves !A(t)$:
\begin{oltableau}
  [\sFmla{\False}{\formula{A}(t)}, just = \TAss
    [\sFmla{\True}{\eq[s][t]}, just = \TAss
      [\sFmla{\True}{\formula{A}(s)}, just = \TAss
        [\sFmla{\True}{\formula{A}(t)}, just={\TRule{\True}{\eq}[2, 3]}, close]
      ]
    ]
  ]
\end{oltableau}
Bu, özdeşlerin birbirinin yerine konabilirliği ilkesi ya da Leibniz Yasası
olarak tanıdık gelebilir.

!!^{tableau}s $\eq$ bağıntısının simetrik olduğunu, yani
$\eq[s_1][s_2] \Proves \eq[s_2][s_1]$ olduğunu ispatlar:
\begin{oltableau}
  [\sFmla{\False}{\eq[s_2][s_1]}, just = \TAss
    [\sFmla{\True}{\eq[s_1][s_2]}, just = \TAss
      [\sFmla{\True}{\eq[s_1][s_1]}, just = {$\eq$}
        [\sFmla{\True}{\eq[s_2][s_1]}, just = {\TRule{\True}{\eq}[2, 3]}, close]
      ]
    ]
  ]
\end{oltableau}
Burada $2$.~satır, $\TRule{\True}{\eq}$ kuralının ilk önkoşul
!!{formula}{}ü~$\sFmla{\True}{\eq[s_1][s_2]}$'dir. $3$.~satır ise
$\sFmla{\True}{!A(s_1)}$ biçimindeki ikinci önkoşuldur---$!A(x)$'i
$\eq[x][s_1]$ olarak düşünün; böylece $!A(s_1)$, $\eq[s_1][s_1]$ ve
$!A(s_2)$, $\eq[s_2][s_1]$ olur.

Ayrıca $\eq$ bağıntısının geçişli olduğunu, yani $\eq[s_1][s_2],
\eq[s_2][s_3] \Proves \eq[s_1][s_3]$ olduğunu da ispatlarlar:
\begin{oltableau}
  [\sFmla{\False}{\eq[s_1][s_3]}, just = \TAss
    [\sFmla{\True}{\eq[s_1][s_2]}, just = \TAss
      [\sFmla{\True}{\eq[s_2][s_3]}, just = \TAss
        [\sFmla{\True}{\eq[s_1][s_3]}, just = {\TRule{\True}{\eq}[3, 2]}, close]
      ]
    ]
  ]
\end{oltableau}
Bu !!{tableau} içinde $\TRule{\True}{\eq}$ kuralının ilk önkoşul
!!{formula}{}ü $3$.~satırdaki $\sFmla{\True}{\eq[s_2][s_3]}$'tür
($s_2$,~$t_1$; $s_3$ ise~$t_2$ rolündedir). $\sFmla{\True}{!A(s_2)}$
biçimindeki ikinci önkoşul $2$.~satırdır. Burada $!A(x)$'i
$\eq[s_1][x]$ olarak düşünün; böylece $!A(s_2)$, $\eq[s_1][s_2]$ (yani
$2$.~satır), $!A(s_3)$ ise sonuçtaki !!{formula}~$\eq[s_1][s_3]$ olur.
\end{ex}

\begin{prob}
Aşağıdakiler için kapalı !!{tableau}s verin:
\begin{enumerate}
\item $\sFmla{\False}{\lforall[x][\lforall[y][((x = y \land !A(x))
      \lif !A(y))]]}$
\item $\sFmla{\False}{\lexists[x][(!A(x) \land
    \lforall[y][(!A(y) \lif y = x)])]}$,\\
  $\sFmla{\True}{\lexists[x][!A(x)] \land
    \lforall[y][\lforall[z][((!A(y) \land !A(z)) \lif y = z)]]}$
\end{enumerate}
\end{prob}

\end{document}
